\documentclass[12pt]{aghdpl}

%---------------------------------------------------------------------------
% Dane do strony tytułowej (wzór WIEiT, ,,Strona tytułowa I stopień'').
% Tytuły muszą się zgadzać co do słowa z tematem zatwierdzonym w USOS.
%---------------------------------------------------------------------------

\author{Norbert Bochenek}
\shortauthor{N. Bochenek}

\titlePL{Wykorzystanie algorytmów uczenia maszynowego oraz modeli typu transformer w~detekcji anomalii w~ruchu sieciowym}
% tytuł angielski wielkimi literami według zasad pisowni tytułów w języku angielskim (wytyczne)
\titleEN{Using Machine Learning Algorithms and Transformer Models to Detect Anomalies in Network Traffic}

% skrócony tytuł do stopki stron
\shorttitlePL{Uczenie maszynowe i transformery w detekcji anomalii sieciowych}
\shorttitleEN{Machine Learning and Transformers in Network Anomaly Detection}

% Wzór WIEiT: ,,Projekt dyplomowy'' na koniec studiów I stopnia, ,,Praca dyplomowa'' na koniec II stopnia.
\thesistype{Projekt dyplomowy}

\supervisor{dr inż. Paweł Kuraś}
\degreeprogramme{Teleinformatyka}
\faculty{Wydział Informatyki, Elektroniki i Telekomunikacji}
\institute{INSTYTUT TELEKOMUNIKACJI I CYBERBEZPIECZEŃSTWA}
\date{2026}

% Podziękowania są opcjonalne (wzór wydziału ich nie wymaga); odkomentuj, żeby dodać stronę z nimi.
% \acknowledgements{Dziękuję \dots}


\begin{document}

	\titlepages
	\RedefinePlainStyle

	% streszczenie i abstract: sekcje nienumerowane (wytyczne)
	\chapter*{Streszczenie}
\thispagestyle{plain}

% Streszczenie po polsku (ok. 1 akapit) i słowa kluczowe.

\noindent\textbf{Słowa kluczowe:} detekcja anomalii, uczenie maszynowe, transformer, ruch sieciowy

\begin{otherlanguage}{english}
\chapter*{Abstract}
\thispagestyle{plain}

% Abstract in English (the same content as above) and keywords.

\noindent\textbf{Keywords:} anomaly detection, machine learning, transformer, network traffic
\end{otherlanguage}


	\setcounter{tocdepth}{2}
	\tableofcontents
	\clearpage

	\chapter{Wstęp}

% Kilka zdań wprowadzenia przed pierwszym podrozdziałem (wytyczne): motywacja (rozdz. 1 w blueprincie w Notion).

\section{Cel pracy}

% Cel pracy i pytania badawcze PB1-PB3.

\section{Zakres pracy}

\section{Struktura pracy}


	\printbibliography

	% spisy rysunków i tabel na końcu pracy
	\listoffigures
	\listoftables

\end{document}
